%Dokumentklasse

%draft als optionohne bilder für bessere performance
%\documentclass[a4paper,12pt,draft]{scrreprt}

%normal mit Bildern
\documentclass[a4paper,12pt]{scrreprt}

\usepackage[left= 3cm,right = 3cm, bottom = 3cm,top = 3cm]{geometry}
%\usepackage[onehalfspacing]{setspace}

% ============= Packages =============
% Dokumentinformationen
\usepackage[
pdftitle={Praktikum - Umwelttechnik},
pdfsubject={},
pdfauthor={Roman-Luca Zank},
pdfkeywords={},	
%Links nicht einrahmen
hidelinks
]{hyperref}

%nur Text zum prüfen des Umfangs

% Standard Packages
%\usepackage[bottom]{footmisc}
\usepackage[utf8]{inputenc}
\usepackage[ngerman]{babel}

\usepackage[T1]{fontenc}
%\usepackage{helvet}

%\renewcommand{\familydefault}{\sfdefault}

\usepackage{graphicx}
\graphicspath{{img/}}
\usepackage{mhchem}
\usepackage{fancyhdr}
\usepackage{lmodern}
\usepackage{color}
\usepackage{placeins}
\usepackage{booktabs}
\usepackage{caption}
\usepackage[list=true]{subcaption}
\usepackage{longtable}
\usepackage{tikz}
\usepackage{pgfplots}
\usepackage{lastpage}

%Einheitenpackage
\usepackage{siunitx}  
\sisetup{	locale = DE, 
	per-mode=fraction,
	inter-unit-product=\ensuremath{\cdot},
	detect-weight = true,
	quotient-mode=fraction
}
%neue Einheiten definieren
\DeclareSIUnit\xyz{xyz}	
\DeclareSIUnit\rpm{rpm}	

%Automatisch cdot statt *
\DeclareMathSymbol{*}{\mathbin}{symbols}{"01}



%\setcounter{lofdepth}{2}      %für subfigures in list of figures

%Tabelle
\usepackage{tabularx}
\usepackage{tabulary}

%%Verzeichnispackages
%\usepackage{natbib}

%\usepackage{hyperref}
%\newcites{

% zusätzliche Schriftzeichen der American Mathematical Society
\usepackage{amsfonts}
\usepackage{amsmath}

%Abkürzungsverzeichnis
\usepackage{acronym}

%kein Abstand bei neuem Kapitel vom Seitenanfang
%\vspace*{2.3\baselineskip} = ORIGINAL
\renewcommand*{\chapterheadstartvskip}{\vspace*{.0\baselineskip}}

%nicht einrücken nach Absatz
\setlength{\parindent}{0pt}

% ============= Kopf- und Fußzeile =============
\pagestyle{fancy}
%
\lhead{}
\chead{}
\rhead{}%\slshape }%\leftmark}
%%
\lfoot{}
\cfoot{}
\rfoot[{\thepage\ of \pageref*{LastPage}}]{Seite \thepage\ von \pageref*{LastPage}}
%%
\renewcommand{\headrulewidth}{0pt}
\renewcommand{\footrulewidth}{0pt}
\renewcommand{\chapterpagestyle}{fancy}

%Fußnotelinie
%\let\footnoterule

%Fußnote mit Klammer
\renewcommand*{\thefootnote}{(\arabic{footnote})}

%Abb. statt Abbildung
\addto\captionsngerman{%
	\renewcommand{\figurename}{Abb.}%
	\renewcommand{\tablename}{Tab.}%
}

% ============= Package Einstellungen & Sonstiges ============= 
%Besondere Trennungen
%\hyphenation{De-zi-mal-tren-nung}
\usepackage[none]{hyphenat}
\hyphenpenalty=5000
\tolerance=5000
\providecommand\phantomsection{}

\usepackage{mathtools}


% ============= Dokumentbeginn =============

\begin{document}
%Seiten ohne Kopf- und Fußzeile sowie Seitenzahl
\pagestyle{empty}

%\begin{center}
\begin{tabular}{p{\textwidth}}


\begin{center}
\includegraphics[scale=0.75]{logos.jpg}\\
\end{center}


\\

\begin{center}
\LARGE{\textsc{
Protokoll \\
Umwelttechnik\\
}}
\end{center}

\\

%\begin{center}
%\large{Fakultät für Muster und Beispiele \\
%der Hochschule Musterhausen \\}
%\end{center}
%
%\\

\begin{center}
\textbf{\Large{V2 - Abwasserbeschaffenheit I}}
\end{center}

\begin{center}
	\large{Gruppe 1.2 (BCUT3)}
\end{center}


\\
%\begin{center}
%zur Erlangung des akademischen Grades\\
%Bachelor of Engineering
%\end{center}


%\begin{center}
%vorgelegt von
%\end{center}

\begin{center}
\Large{\textbf{Teilnehmer:}} \\ 
\end{center}
\begin{center}
\large{Christoph Hecht \\
	Willy Messerschmidt \\
	Roman-Luca Zank} \\
\end{center}


\\

\begin{center}
\begin{tabular}{lll}
\large{\textbf{Protokollführer:}} & & \large{Willy Messerschmidt} \\
& & \href{mailto:willy.messerschmidt@stud.hs-merseburg.de}{{\footnotesize willy.messerschmidt@stud.hs-merseburg.de}}\\
&&\\
\large{\textbf{Datum der Versuchsdurchführung:}}&& \large{03.12.2019}\\
&&\\
\large{\textbf{Abgabedatum:}}&& \large{17.12.2019}
\end{tabular}
\end{center}

\\ \\ \\ \\ \\ \\
\large{Merseburg den \today}

\end{tabular}
\end{center}


%\include{14_danksagungen}

%\include{15_zusammenfassung}

% Beendet eine Seite und erzwingt auf den nachfolgenden Seiten die Ausgabe aller Gleitobjekte (z.B. Abbildungen), die bislang definiert, aber noch nicht ausgegeben wurden. Dieser Befehl fügt, falls nötig, eine leere Seite ein, sodaß die nächste Seite nach den Gleitobjekten eine ungerade Seitennummer hat. 
\cleardoubleoddpage

% Pagestyle für Titelblatt leer
\pagestyle{empty}

%Seite zählen ab
\setcounter{page}{0}

%Titelblatt
\begin{center}
\begin{tabular}{p{\textwidth}}


\begin{center}
\includegraphics[scale=0.75]{logos.jpg}\\
\end{center}


\\

\begin{center}
\LARGE{\textsc{
Protokoll \\
Umwelttechnik\\
}}
\end{center}

\\

%\begin{center}
%\large{Fakultät für Muster und Beispiele \\
%der Hochschule Musterhausen \\}
%\end{center}
%
%\\

\begin{center}
\textbf{\Large{V2 - Abwasserbeschaffenheit I}}
\end{center}

\begin{center}
	\large{Gruppe 1.2 (BCUT3)}
\end{center}


\\
%\begin{center}
%zur Erlangung des akademischen Grades\\
%Bachelor of Engineering
%\end{center}


%\begin{center}
%vorgelegt von
%\end{center}

\begin{center}
\Large{\textbf{Teilnehmer:}} \\ 
\end{center}
\begin{center}
\large{Christoph Hecht \\
	Willy Messerschmidt \\
	Roman-Luca Zank} \\
\end{center}


\\

\begin{center}
\begin{tabular}{lll}
\large{\textbf{Protokollführer:}} & & \large{Willy Messerschmidt} \\
& & \href{mailto:willy.messerschmidt@stud.hs-merseburg.de}{{\footnotesize willy.messerschmidt@stud.hs-merseburg.de}}\\
&&\\
\large{\textbf{Datum der Versuchsdurchführung:}}&& \large{03.12.2019}\\
&&\\
\large{\textbf{Abgabedatum:}}&& \large{17.12.2019}
\end{tabular}
\end{center}

\\ \\ \\ \\ \\ \\
\large{Merseburg den \today}

\end{tabular}
\end{center}
 %Prokolle
%\begin{center}
\begin{tabular}{p{\textwidth}}


\begin{center}
\includegraphics[scale=0.75]{img/logos.jpg}\\
\end{center}


\\

\begin{center}
\LARGE{\textsc{
Recherche \\
Rückgewinnung von Ammoniak aus Industrieabwässern\\
}}
\end{center}

%\begin{center}
%\large{Fakultät für Muster und Beispiele \\
%der Hochschule Musterhausen \\}
%\end{center}
%
%\\
 \\
 
\begin{center}
\textbf{\Large{Seminararbeit in Medienrecherche}}
\end{center}

\begin{center}
	\large{im WiSe 2019}
\end{center}
 \\
%\begin{center}
%zur Erlangung des akademischen Grades\\
%Bachelor of Engineering
%\end{center}


\begin{center}
\large{vorgelegt von}
\end{center}
\\


\begin{center}
\Large{\textbf{Roman-Luca Zank}} \\
\end{center}

\begin{center}
3. Semester \\
Chemie- und Umwelttechnik \\
\end{center}


\begin{center}
\begin{tabular}{lll}
	\textbf{E-Mail:} & & romanzank@mail.de\\
	\textbf{Matrikelnummer:} & &25240\\
	\textbf{Adresse:} & &Platz der Bausoldaten 2, Zimmer 224\\
	\textbf{Ort:} & &06217 Merseburg\\
	&& \\
	\textbf{Prüfer:} & & Dr. Frank  Baumann\\
\end{tabular}
\end{center}

\\ \\ \\ \\ \\
\large{Merseburg, \today}

\end{tabular}
\end{center}
 %Seminar-/Abschlussarbeit

% Pagestyle für Rest des Dokuments
\pagestyle{fancy}

%Inhaltsverzeichnis
\tableofcontents
\thispagestyle{empty}

%Inhalt

%%Verzeichnis aller Bilder
\label{sec:bilder}
\listoffigures
\addcontentsline{toc}{chapter}{Abbildungsverzeichnis}
\thispagestyle{empty}

%Verzeichnis aller Tabellen
\label{sec:tabellen}
\listoftables
\addcontentsline{toc}{chapter}{Tabellenverzeichnis}
\thispagestyle{empty}

%Abkürzungsverzeichnis
\setlength{\columnsep}{20pt}
%\twocolumn
\addchap{Nomenklatur}
\label{sec:abkurzung}
\begin{acronym}
\acro{kf}[$\text{k}_\text{f}$]{Durchlässigkeitsbeiwert}
\acro{t}{Durchlaufzeit}
\acro{tm}[$\text{t}_\text{m}$]{Mittlere Durchlaufzeit}
\acro{V}{Volumen}
\acro{h}{Höhe der Wassersäule}
\acro{Q}{Volumenstrom}
\acro{l}{Durchströmte Länge}
\acro{A}{Grundfläche}
\acro{d}{Durchmesser}

\end{acronym}
%\subsubsection{Aufrufen einer Abkürzung}
%\acs{rT}
%\begin{verbatim}
%\acs{Abkürzung}
%\end{verbatim}

\chapter{Aufgabenstellung}
\label{sec:aufgabenstellung}
%In der Aufgabenstellung wird (in eigenen Worten und ganzen Sätzen) formuliert, was das Ziel des 
%Versuches ist.  
%[Beachten Sie die eigentliche Aufgabenstellung in den Versuchsanleitungen sowie die Hinweise zur Auswertung!] 

Im Versuch 4 "`Bodencharakterisierung - Durchlässigkeitsbeiwert"' wird die Wasserdurchlässigkeit verschiedener Böden mit unterschiedlichen Korngrößen untersucht. \\
Die ermittelten Ergebnisse sind in grafischer, sowie in tabellarischer Form darzustellen. Ziel ist es die Bodenart anhand des $\text{k}_\text{f}$-Wertes nach der DIN 18130-1 zu bestimmen.


\chapter{Geräte und Chemikalien}
\label{sec:geraete}

\textbf{Geräte:}
\begin{itemize}
	\item Magnetrührer mit Rührfisch
	\item Bechergläser
	\item Erlenmeyerkolben
	\item Büchnertrichter für Vakuumfiltration
	\item Erlenmeyerkolben mit angeschlossener Saugleitung
	\item Filterpapier
	\item Schnelltests von \textsc{Chemsolute$^{\textsuperscript{\textregistered}}$}:
		\begin{itemize}
			\item j
		\end{itemize}
	\item Reflectoquant$^{\textsuperscript{\textregistered}}$ für Reflektometer:
		\begin{itemize}
			\item j
		\end{itemize}
	\item Reflektometer RQflex$^{\textsuperscript{\textregistered}}$ plus 10 von \textsc{Merck}
\end{itemize}

\vspace*{5mm}

\textbf{Proben/Chemikalien:}
\begin{itemize}
	\item destilliertes Wasser
	\item Abwasserproben 1, 2 \& 3
\end{itemize}







\chapter{Durchführung}
\label{sec:durchfuerung}



\chapter{Ergebnisse}
\label{sec:ergebnisse}

Im folgenden Protokollabschnitt werden die Versuchsergebnisse der Versuchsdurchführung präsentiert.\\

Die aus Versuchsteil 1 wahrgenommenen Einstufungen der Abwasserproben 1 bis 3 sind in Tab. \ref{tab:optisch_sp} aufgeführt und im Vergleich mit Abb. \ref{fig:proben} nachzuvollziehen.
\vspace*{-2.5mm}
\renewcommand{\arraystretch}{1.2}
\begin{table}[h!]
	\centering
	\caption{Wahrgenommene Einstufungen der Färbung und Trübung der Abwasserproben 1 bis 3}
	\label{tab:optisch_sp}
	%\resizebox{10cm}{!}{
	\begin{tabulary}{\textwidth}{C|CCC}
		\hline
		\textbf{} & \textbf{Probe 1} & \textbf{Probe 2} & \textbf{Probe 3} \\ 
		\hline
		\textbf{Färbung} & stark gefärbt& stark gefärbt& schwach \\
		\textbf{Trübung} & undurchsichtig & undurchsichtig& klar \\
		\hline
	\end{tabulary}
	%}
\end{table}
\FloatBarrier
\vspace*{-2.5mm}

%Tabelle Ende

Für die Abwasser Proben 1 bis 3 sind die wahrgenommenen Einschätzungen zu den Gerüchen in Tab. \ref{tab:olfaktorisch_sp} aufgeführt.

\vspace*{-2.5mm}
\renewcommand{\arraystretch}{1.2}
\begin{table}[h!]
	\centering
	\caption{Wahrgenommene Einstufungen des Geruchs der Abwasserproben 1 bis 3}
	\label{tab:olfaktorisch_sp}
	%\resizebox{10cm}{!}{
	\begin{tabulary}{\textwidth}{C|CCC}
		\hline
		\textbf{} & \textbf{Probe 1} &  \textbf{Probe 2}&  \textbf{Probe 3}\\
		\hline
		\textbf{Intensität}	& schwach	&stark &ohne\\
		\textbf{Art}		& erdig, muffig	& modrig	&-\\
		\textbf{charak.-chemisch  }	& - & - & -\\
		\hline
	\end{tabulary}
	%}
\end{table}
\FloatBarrier
\vspace*{-0.5mm}

%Tabelle Ende

Die elektrochemischen Messwerte für die Abwasserproben 1 bis 3 sind in Tab. \ref{tab:elektrochemisch} eingetragen.
Zu beachten ist dabei, dass die Messwerte für den pH-Wert nur eine grobe Einschätzung liefern, da laut Messgerät die pH-Wert-Elektrode hätte gewechselt werden müssen. \\
Weitere Diskussionen zur Fehlerbetrachtung sind unter Abschnitt \ref{sec:fehler} zu finden.

\vspace*{-2.5mm}
\renewcommand{\arraystretch}{1.2}
\begin{table}[h!]
	\centering
	\caption{Elektrochemische Messwerte der Abwasserproben 1 bis 3}
	\label{tab:elektrochemisch}
	%\resizebox{10cm}{!}{
	\begin{tabulary}{\textwidth}{L|C|C|C}
		\hline
		\textbf{} 						& \textbf{pH-Wert\protect\footnotemark[1] $\boldsymbol{\left[-\right]}$} & \textbf{Leitfähigkeit $\boldsymbol{\left[\si{\milli \siemens \per \centi \meter}\right]}$}& \textbf{Temperatur} $\boldsymbol{\left[\si{\celsius}\right]}$\\
		\hline
		\textbf{Probe 1}				& 1,70	& 10,59	& 17,5	\\
		\textbf{Probe 1 (filtr.)}		& 1,68  & 10,32	& 17,6	\\
		\hline
		\textbf{Probe 2}				& 7,38 	& 2,67	& 17,4	\\
		\textbf{Probe 2 (filtr.)}		& 7,58	& 2,55	& 17,8	\\
		\hline
		\textbf{Probe 3}				& 7,44	& 2,43	& 17,4	\\
		\textbf{Probe 3 (filtr.)}		& 7,54	& 2,27	& 17,9	\\
		\hline
	\end{tabulary}
	%}
\end{table}
\FloatBarrier
\vspace*{-2.5mm}

%Tabelle Ende


\footnotetext[1]{pH-Wert-Elektrode auf \SI{25}{\celsius} kalibriert}


\newpage 

Die für die fotometrischen bzw. kolorimetrischen Messungen erhobenen Daten lassen sich für die Abwasserproben 1 bis 3 in Tab. \ref{tab:fotometrisch} wiederfinden. In dieser Tabelle sind sowohl die Schnelltestergebnisse (ST) als auch die Reflektometerergebnisse (RM) aufgeführt. \\
Ebenfalls aufgeführt ist die Berechnung der Konzentration an Nitrat \ce{NO3-} mittels Schnelltest für Probe 1, da der Messbereich von maximalen \SI{500}{\milli \gram \per \liter} nicht ausreicht. Die dafür nötige Berechnung der Verdünnung der Probe 1 ist ab Gleichung \ref{gl1} gezeigt.

\vspace*{-2.5mm}
\renewcommand{\arraystretch}{1.2}
\begin{table}[h!]
	\centering
	\caption{Foto-/kolorimetrisch Messwerte der Abwasserproben 1 bis 3}
	\label{tab:fotometrisch}
	%\resizebox{10cm}{!}{
	\begin{tabulary}{1.2\textwidth}{L|C|C|C}
		\hline
		\textbf{} 						& \textbf{Probe 1} $\boldsymbol{\left[\si{\milli \gram \per \liter}\right]}$& \textbf{Probe 2}$\boldsymbol{\left[\si{\milli \gram \per \liter}\right]}$&\textbf{Probe 3}$\boldsymbol{\left[\si{\milli \gram \per \liter}\right]}$\\
		\hline
		\textbf{Nitrit \ce{NO2-} (Schnelltest)}	& 2	& 0&2\\
		\textbf{Nitrit \ce{NO2-} (Reflektometer)} & <0,3  & <0,3	&<0,3\\
		\hline
		\textbf{Nitrat \ce{NO3-} (ST)}			& >500 	& 0	&0\\
		\textbf{Nitrat \ce{NO3-} (RM)}	{\footnotesize (verdünnt/unverdünnt)}		& 48/480	& <3	&7,4\\
		\hline
		\textbf{Phosphat \ce{PO4^3-} (ST)}		& {\footnotesize kein Messwert}\protect\footnotemark[2]	& 0	&25\\
		\textbf{Phosphat \ce{PO4^3-} (RM)}		& {\footnotesize kein Messwert}\protect\footnotemark[2]	& 12	&27\\
		\hline
		\textbf{Ammonium \ce{NH4+} (RM)}		& {\footnotesize kein Messwert}\protect\footnotemark[2]		&	1,3	&0,6\\
		\hline
	\end{tabulary}
	%}
\end{table}
\FloatBarrier
\vspace*{-2.5mm}

%Tabelle Ende

\footnotetext[2]{keine Messung möglich, da laut Packungsbeilage  der pH-Wert zu niedrig für ein verwertbares Ergebnis ist}

\subsubsection{\underline{Verdünnung für Nitrat-Gehaltbestimmung der Probe 1}}
\textbf{Es wird von $\boldsymbol{V_{\text{Probe 1}} = \SI{100}{\milli \liter}}$ ausgegangen:}
\begin{flalign}
\label{gl1}
	\frac{\SI{500}{\milli \gram \per \liter}}{\SI{50}{\milli \gram \per \liter}} &= \frac{\SI{100}{\milli \liter}}{V_{\text{zu verdünnen}}}\\[2mm]
	V_{\text{zu verdünnen}}	&= \SI{100}{\milli \liter}*\frac{\SI{50}{\milli \gram\per \liter}}{\SI{500}{\milli \gram\per \liter}}\\[-2mm]
			&=\underline{ \SI{10}{\milli \liter}} \\[3mm]
			V_{\ce{H2O}}	&= \SI{100}{\milli \liter}-V_{\text{zu verdünnen}}\\
					&=\SI{100}{\milli \liter}-\SI{10}{\milli \liter}\\
					&= \underline{\SI{90}{\milli \liter}}\\[3mm]
	\frac{\SI{10}{\milli \liter}}{\SI{100}{\milli \liter}}	&=\frac{1}{10}\\
					&=\underline{\underline{1:10}} \quad \text{{\footnotesize (Verdünnung)}}
\end{flalign}

Mit \SI{50}{\milli \gram \per \liter} ist der Messbereich des Schnellteststreifens für Nitrat mit \SI{0}{}-\SI{500}{\milli \gram \per \liter} gut erfasst und das Volumen von \SI{10}{\milli \liter} ist einfach mittels Messzylinder abzumessen.\\
Daraus folgt eine Verdünnung von 1:10 mit einem Volumenteil (\SI{10}{\milli \liter}) Abwasserprobe 1 und neun Volumenteilen (\SI{90}{\milli \liter}) destilliertes Wasser.\\

\newpage

Der Messwert des Reflektometers für die verdünnte Lösung der Probe 1 ergibt für den Nitratgehalt \SI{48}{\milli \gram \per \liter} (siehe Tab. \ref{tab:elektrochemisch}). Hochgerechnet auf die unverdünnte Lösung ergibt sich:
\begin{flalign}
	\frac{c_{1,\text{verdünnt}}}{c_{1,\text{unverdünnt}}}	&= \frac{1}{10}\\
	c_{1,\text{unverdünnt}} &= 10*c_{1,\text{verdünnt}}\\
						&= 10*\SI{48}{\milli \gram \per \liter}\\
						&=\underline{\underline{ \SI{480}{\milli \gram \per \liter}}}
\end{flalign} 

\subsubsection{\underline{Bestimmung des anorganisch gebundenen Gesamtstickstoffs}}
Der gesamte anorganisch gebundene Stickstoff kann über den Masseanteil des Elementaren Stickstoffs an den einfachen anorganischen Verbindungen Nitrat-, Nitrit- und Ammoniumionen berechnet werden.\\
Dabei kann aus dem Massenanteil des Stickstoffs am Molekül abgeleitet werden, dass $\SI{1}{\milli\gram\per\liter}$Ammoniumionen $\frac{7}{9}\si{\milli\gram\per\liter}$ anorganisch gebundenem Stickstoff entsprechen.(vgl. Gleichung 5.1)
\begin{flalign}
\frac{M(N)}{M(\ce{NO_3^-})} &= \frac{\SI{14}{\gram\per\mole}}{\SI{14}{\gram\per\mole}+3*\SI{16}{\gram\per\mole}} =\frac{7}{31} \approx 22,6\% \\
\frac{M(N)}{M(\ce{NO_2^-})} &= \frac{\SI{14}{\gram\per\mole}}{\SI{14}{\gram\per\mole}+2*\SI{16}{\gram\per\mole}} =\frac{7}{23}\approx 30,4\%\\
\frac{M(N)}{M(\ce{NH_4^+})} &= \frac{\SI{14}{\gram\per\mole}}{\SI{14}{\gram\per\mole}+4*\SI{1}{\gram\per\mole}} =\frac{7}{9} \approx 70,0\% 
\end{flalign}

Zur Berechnung des anorganisch gebundenen Gesamtstickstoffs können die Ionenkonzentrationen mit dem entsprechenden Umrechnungsfaktoren aufsummiert werden. Beispielhaft wird $N_{ges}$ für Probe 2 mit de Reflektometerergebnissen vorgerechnet, zu sehen in den Gleichungen \ref{gl:tin1} bis \ref{gl:tin2}. Die restlichen Werte des $N_{ges}$ sind in Tab. \ref{tab:tin} dargestellt.
\begin{flalign}
\label{gl:tin1}
	N_{ges,i} 	&= c_{(i,\ce{NO3^-})}*22,6\%+c_{(i,\ce{NO2^-})}*30,4\%+c_{(i,\ce{NH4^+})}*70,0\%\\[2mm]
			&= \SI{7,4}{\milli\gram\per\liter}*22,6\%+\SI{0,3}{\milli\gram\per\liter}*30,4\%+\SI{0,6}{\milli\gram\per\liter}*70,0\%\\
			&= \SI{1,67}{\milli\gram\per\liter}+\SI{0,09}{\milli\gram\per\liter}+\SI{0,42}{\milli\gram\per\liter}\\
\label{gl:tin2}
	N_{ges,3}	&=\underline{\underline{\SI{2,18}{\milli\gram\per\liter}}}
\end{flalign}

\vspace*{-8.5mm}
\renewcommand{\arraystretch}{1.2}
\begin{table}[h!]
	\centering
	\caption{$N_{ges}$ der Abwasserproben 1 bis 3}
	\label{tab:tin}
	%\resizebox{10cm}{!}{
	\begin{tabulary}{1.2\textwidth}{C|C|C|C}
		\hline
		\textbf{} 						& \textbf{Probe 1} & \textbf{Probe 2}&\textbf{Probe 3}\\
		\hline
		\textbf{$N_{ges}$ $\boldsymbol{\left[\si{\milli \gram \per \liter}\right]}$}	& 108,67\protect\footnotemark[3] & 0,77 & 2,18\\
		\hline
	\end{tabulary}
	%}
\end{table}
\FloatBarrier
\vspace*{-2.5mm}

%Tabelle Ende
\footnotetext[3]{zweifelhaftes Ergebnis, da kein Messwert für den Ammoniumgehalt vorhanden ist}

\newpage

\subsection*{\underline{Bestimmung des Phosphorgehalts}}

\begin{flalign}
\frac{M(P)}{M(\ce{PO_4^3-})} &= \frac{\SI{31}{\gram\per\mole}}{\SI{31}{\gram\per\mole}+4*\SI{16}{\gram\per\mole}} =\frac{31}{95} \approx 32,6\% 
\end{flalign}

Beispielhaft für Reflektometerergebnisse der Probe 3 vorgerechnet:
\begin{flalign}
P_{ges,i}	&= c_{(i,\ce{PO4^3-})}*32,6\% \\[2mm]
P_{ges,3}	&= \SI{27}{\milli \gram \per \liter}*32,6\%\\[2mm]
		&= \underline{\underline{\SI{8,8}{\milli \gram \per \liter}}}
\end{flalign}

\vspace*{-3.5mm}
\renewcommand{\arraystretch}{1.2}
\begin{table}[h!]
	\centering
	\caption{$P_{ges}$ der Abwasserproben 1 bis 3}
	\label{tab:tp}
	%\resizebox{10cm}{!}{
	\begin{tabulary}{1.2\textwidth}{C|C|C|C}
		\hline
		\textbf{} 						& \textbf{Probe 1} & \textbf{Probe 2}&\textbf{Probe 3}\\
		\hline
		\textbf{$P_{ges}$ $\boldsymbol{\left[\si{\milli \gram \per \liter}\right]}$}	& {\footnotesize kein Ergebnis\protect\footnotemark[4]} & 3,9 & 8,8\\
		\hline
	\end{tabulary}
	%}
\end{table}
\FloatBarrier
\vspace*{-3.5mm}
%Tabelle Ende
\footnotetext[4]{kein Ergebnis, da kein Messwert für den Phosphatgehalt vorhanden ist}

%\newpage

\subsection*{Gegenüberstellung der Mindestanforderungen für das Einleiten kommunaler Abwässer in den Vorfluter der GK 5 mit den Abwasserproben}
Die Referenzwerte der Mindestanforderungen für das Einleiten kommunaler Abwässer in den Vorfluter der Größenklasse 5 ist im Anhang von \cite[S. 29]{Skript} zu finden.
\begin{figure}[h!]
	\begin{tikzpicture}
	\begin{axis}[
	xbar=1pt,% space of 0pt between adjacent bars
	bar width=7,
	width=15cm,
	height=7cm,
	%minor y tick num=4,
	xmax=110,xmin=0,
	x tick label style={/pgf/number format/.cd,%
		scaled x ticks = false,
		set decimal separator={,},
		fixed},
	symbolic y coords={Probe 3,Probe 2,Probe 1,max. für GK 5},
	ytick=data,
	%legend style={at={(0.5,-0.15)},
	%	anchor=north,legend columns=-1},
	xtick={0,10,...,110},
	grid=major,
	xlabel=Gehalt in  \si{\milli\gram\per\liter},
	%enlargelimits=0.15,
	]
	%NH4-N
	\addplot coordinates {
		(0,Probe 1) (1.3,Probe 2) (0.6,Probe 3) (10,max. für GK 5)
	};
	%N gesamt
	\addplot coordinates {
		(108,Probe 1) (0.77,Probe 2) (2.18,Probe 3) (13,max. für GK 5)
	};
	%P gesamt
	\addplot coordinates {
		(0,Probe 1) (3.9,Probe 2) (8.8,Probe 3) (1,max. für GK 5) 
	};
	\legend{$NH_4-N$,$N_{ges}$,$P_{ges}$}
	\end{axis}
	\end{tikzpicture}
	\caption{Gegenüberstellung der Mindestanforderungen für das Einleiten kommunaler Abwässer in den Vorfluter der GK 5 \linebreak für die Abwasserproben 1 bis 3}
	\label{Balkendiagramm}
\end{figure}
\FloatBarrier

\subsubsection{Gegenüberstellung der Messwerte mit der durchschnittlichen Beschaffenheit von häuslichem Abwasser}

Um die Messwerte des Versuches mit häuslichem Abwasser gegenüberzustellen wird die Tabelle Tab. \ref{label} (siehe \cite[S. 29]{Skript}) genutzt.

%\vspace*{-2.5mm}
\renewcommand{\arraystretch}{1.2}
\begin{table}[h!]
	\centering
	\caption{Tabellarischer Vergleich der Messwerte mit den Mindestanforderungen für das Einleiten kommunaler Abwässer in den Vorfluter der GK 5}
	\label{tab_vgl}
	%\resizebox{10cm}{!}{
	\begin{tabulary}{1.2\textwidth}{|l|C|C|C|}
		\hline
		Werte in \si{\milli\gram\per\liter} & \textbf{$NH_4^+-N$} & \textbf{$N_{ges}$}&\textbf{$P_{ges}$}\\
		\hline
		\textbf{Grenzwert} & \textbf{10} & \textbf{13} & \textbf{1} \\
		\hline
		Probe 1 & 0 & 108 & 0 \\
		Probe 2 & 1,3 & 0,77 & 3,9 \\
		Probe 3 & 0,6 & 2,18 & 8,8 \\
		\hline
	\end{tabulary}
	%}
\end{table}
\FloatBarrier

\begin{figure}[h!]
	\begin{tikzpicture}
	\begin{axis}[
	%x tick label style={
	%	/pgf/number format/1000 sep=},
	xtick = data,
	ylabel=pH-Wert,
   	enlarge x limits=0.15,
	ybar,
	ymin = 0,
	ymax = 14,
	bar width=10pt,
	width=15cm,
	height=9cm,
	symbolic x coords={0,Probe 1, Probe 2,Probe 3,2,1},
	]
	\addplot coordinates {(Probe 1,1.7) (Probe 2,7.38) (Probe 3,7.44)};
	\addplot[green,sharp plot,update limits=false, dashed] 
	coordinates {(0,6.6) (1,6.6)} 
	node[above] at (axis cs:Probe 2,6.6) {geringe Belastung};
	\addplot[yellow,sharp plot,update limits=false, dashed] 
	coordinates {(0,7.6) (1,7.6)} 
	node[above] at (axis cs:Probe 2,7.6) {mittelere Belastung};
	\addplot[red,sharp plot,update limits=false, dashed] 
	coordinates {(0,8.6) (1,8.6)} 
	node[above] at (axis cs:Probe 2,8.6) {starke Belastung};
	\legend{pH-Wert}
	\end{axis}
	\end{tikzpicture}
	\caption{pH-Wert-Beschaffenheit der Abwasserproben 1 bis 3}
\end{figure}
\FloatBarrier

\begin{figure}[h!]
	\begin{tikzpicture}
	\begin{axis}[
	%x tick label style={
	%	/pgf/number format/1000 sep=},
	xtick = data,
	ylabel=$N_{ges}$ in \si{\milli \gram \per \liter},
	enlarge x limits=0.15,
	ybar,
	ymin = 0,
	ymax = 120,
	bar width=10pt,
	width=15cm,
	height=9cm,
	symbolic x coords={0,Probe 1, Probe 2,Probe 3,2,1},
	]
	\addplot coordinates {(Probe 1,108.7) (Probe 2,0.77) (Probe 3,2.18)};
	\addplot[green,sharp plot,update limits=false, dashed] 
	coordinates {(0,20) (1,20)} 
	node[above] at (axis cs:Probe 2,20) {geringe Belastung};
	\addplot[yellow,sharp plot,update limits=false, dashed] 
	coordinates {(0,50) (1,50)} 
	node[above] at (axis cs:Probe 2,50) {mittelere Belastung};
	\addplot[red,sharp plot,update limits=false, dashed] 
	coordinates {(0,85) (1,85)} 
	node[above] at (axis cs:Probe 2,85) {starke Belastung};
	\legend{$N_{ges}$-Wert}
	\end{axis}
	\end{tikzpicture}
	\caption{$N_{ges}$-Beschaffenheit der Abwasserproben 1 bis 3}
\end{figure}
\FloatBarrier

\begin{figure}[h!]
	\begin{tikzpicture}
	\begin{axis}[
	%x tick label style={
	%	/pgf/number format/1000 sep=},
	xtick = data,
	ylabel=$P_{ges}$ in \si{\milli \gram \per \liter},
	enlarge x limits=0.15,
	ybar,
	ymin = 0,
	ymax = 50,
	bar width=10pt,
	width=15cm,
	height=9cm,
	symbolic x coords={0,Probe 1, Probe 2,Probe 3,2,1},
	]
	\addplot coordinates {(Probe 1,0) (Probe 2,3.9) (Probe 3,8.8)};
	\addplot[green,sharp plot,update limits=false, dashed] 
	coordinates {(0,7) (1,7)} 
	node[above] at (axis cs:Probe 2,7) {geringe Belastung};
	\addplot[yellow,sharp plot,update limits=false, dashed] 
	coordinates {(0,20) (1,20)} 
	node[above] at (axis cs:Probe 2,20) {mittelere Belastung};
	\addplot[red,sharp plot,update limits=false, dashed] 
	coordinates {(0,30) (1,30)} 
	node[above] at (axis cs:Probe 2,30) {starke Belastung};
	\legend{$P_{ges}$-Wert}
	\end{axis}
	\end{tikzpicture}
	\caption{$P_{ges}$-Beschaffenheit der Abwasserproben 1 bis 3}
\end{figure}
\FloatBarrier

Bli bla blub

\chapter{Diskussion}
\label{sec:diskussion}
In diesem Abschnitt des Protokolls werden nun die Ergebnisse des Abschnittes \ref{sec:ergebnisse} diskutiert und ausgewertet.\\\\
Begonnen wird mit der Auswertung der Sinnesprüfung und der Einschätzung über die Herkunft der Abwasserproben 1 bis 3.\\ 
Die Wasserprobe 1 ist sehr trüb und riecht leicht erdig und muffig. Die starke Trübung und die dunkelbraune bis schwarze Farbe ist wahrscheinlich organischen Ursprungs. Es könnte sich hierbei um Wasser aus einem Sumpfgebiet handeln. Es kann kein Wasser sein, dass bereits durch den Boden gefiltert wurde.
Die für Moore typische anaerobe Zersetzung führt zu einem absinken des pH-Wertes. Der gleiche Prozess könnte auch in einem Faulturm stattgefunden haben. Tatsächlich wurde ein sehr niedriger pH-Wert für diese Probe von 1,7 festgestellt (siehe Tab \ref{tab:elektrochemisch}). Der niedrige pH-Wert verhinderte darauffolgende Phosphat- und Ammoniummessungen mittels foto- bzw. kolorimetrischer Messmethoden.\\ \\
Die Wasserprobe 2 ist ebenfalls sehr dunkel gefärbt und sehr trüb. Noch dazu riecht sie stark moderig und faulig. Es handelt sich hierbei wahrscheinlich um ein Abwasser, in welchem organische Anteile von Mikroorganismen unter Ausstoß von Schwefelwasserstoff und anderen charakteristischen Geruchsstoffen zersetzt werden oder wurden. Auch dieses Wasser wird höchstwahrscheinlich von der Oberfläche stammen. Die Probe enthält im Vergleich sehr wenig Ammonium- als auch Phosphat- und Nitrationen (siehe Tab. \ref{tab:elektrochemisch}). \textit{Diese müssen noch in der Biomasse gebunden vorliegen.}\\ \\
Die Wasserprobe 3 ist nur schwach getrübt und es ist nur ein sehr schwacher Geruch wahrnehmbar. Es sind vergleichsweise hohe Konzentrationen an Phosphat- und Nitrationen nachweisbar (siehe Tab. \ref{tab:elektrochemisch}). Gegenüber der Probe 2 liegen weniger Ammoniumionen vor. Es liegt nah, dass die Probe 3 bereits einen Klärprozess durchlaufen hat, bei dem organische Trübstoffe abgebaut und die enthaltenen Nährstoffe freigesetzt wurden sein könnten.\\ \\
In Anbetracht dessen, dass die Proben im Rahmen eines Praktikums gegeben sind, wäre es möglich, dass sie aus verschieden Phasen der Abwasserbehandlung stammen. In diesem Fall könnte die Probe 2 das ankommende Schmutzwasser, die Probe 2 das Wasser während der Belüftung und die Probe 3 Wasser vor der Phosphatabscheidung darstellen. \\
\newpage
Das Hauptaugenmerk der Untersuchungen liegt auf Phosphat- und Stickstoffverbindungen, da diese wichtige und in vielen Fällen limitierende Nährstoffe für das Pflanzenwachstum darstellen.
Stickstoffverbindungen gelangen dabei als Abbauprodukte organischer Verbindungen wie Harnstoffe oder Fäkalien, sowie aus industriellen Abwässern an. Gerade die Umwandlung von Ammonium- in Nitratverbindungen verbraucht Sauerstoff, der beispielsweise dem Ökosystem See fehlt und zum Fischsterben führen kann. Zudem verschärft sich dieser Sauerstoffverbrauch immer weiter, da die entstehenden Nitrate als Dünger für Pflanzen wirken und bei deren Absterben ebenfalls wieder Sauerstoff für die Zersetzung nötig ist \cite{NundP}. Des Weiteren kann Ammoniak gelöst bereits in geringen Konzentrationen zu einem Sterben von Wasserlebewesen führen \cite{AnalyseAbwasserN}. \linebreak
Phosphorverbindungen wirken ähnlich sauerstoffzehrend wie Nitratverbindungen durch ihren Düngungseffekt bei Pflanzen. Diese Verbindungen stammen meist aus anorganischen Düngemitteln, welche in der Landwirtschaft eingesetzt werden und über Versickerung in das Grundwasser und somit in weitere Gewässer gelangt \cite{NundP,AnalyseAbwasserP}. \linebreak
Somit bestimmen die Konzentrationen der genannten Verbindungen das Algen- bzw. Pflanzenwachstum in den Gewässern in besonderem Maße und sind damit für das Umkippen von aquatischen Ökosystemen verantwortlich. Um die Gewässer vor solch einer sogenannten Eutrophierung zu bewahren, müssen die Abwässer von obigen Nährstoffen befreit werden.\\\\
Im Weiteren folgt die Diskussion, ob die elektrochemischen, fotometrischen und kolorimetrischen Bestimmungen der Abwasserproben den Mindestanforderungen für das Einleiten kommunaler Abwässer in einen Vorfluter der Größenklasse 5 entsprechen \cite[S. 29]{Skript}.\\
Die Mindestanforderungen für das Einleiten kommunaler Abwässer in den Vorfluter bei einer Anlage der Größenklasse 5 sind bei den getesteten Wasserproben 1 bis 3 nur teilweise erfüllt, wie in Tab. \ref{tab_vgl} und Abb. \ref{Balkendiagramm} nachzuvollziehen. Zu erkennen ist, dass keine der drei Proben den Grenzwert für den Ammoniumgehalt von \SI{10}{\milli \gram \per \liter} überschreitet. Stark auffallend ist dafür jedoch die Konzentration des anorganisch gebundenen Gesamtstickstoffs, welche durch die Probe 1 um mehr als das achtfache überschritten wird. Dieses Abwasser müsste zur Einhaltung des Grenzwertes weiter denitrifiziert werden. Die Proben 2 und 3 halten den Grenzwert für anorganisch gebundenen Gesamtstickstoff ein. Den Abwässern der Proben 2 und 3 muss Phosphor entzogen werden, da sie den Grenzwert von \SI{1}{\milli\gram\per\liter} beinahe um das vier- bzw. achtfache überschreiten. Für Probe 1 ist aufgrund des zu geringen pH-Wertes keine Bestimmung des Phosphatgehaltes mit gegebenen Mitteln möglich. Eine Belastung der Probe 1 mit Phosphaten ist daher nicht auszuschließen und müsste über alternative Messmethoden bestimmt werden.\\
Die Abbildungen \ref{ph} bis \ref{pges} stellen die gemessenen Parameter der Abwasserproben im Vergleich zu üblichen häuslichen Abwässern dar \cite[S. 29]{Skript}. Aus der Abbildung \ref{ph} in Bezug auf den pH-Wert geht hervor, dass die Proben 2 und 3 nahezu mittelstark belastet sind. Der etwas basische pH-Wert ließe sich zum Beispiel sehr kostengünstig durch die Zugabe kleiner Mengen einer verträglichen Säure korrigieren. Die Einstufung der Belastungsgrade anhand des pH-Wertes in der Art wie es hier erfolgen soll, weist eine entscheidende Schwäche auf:\linebreak
Abwässer welche sehr sauer sind, gelten nach Tab. \ref{tab:komm} als gering belastet. Ein Abwasser mit dem pH-Wert von Probe 1 kann in natürlichen Gewässern und an der Abwasserentsorgungsinfrastruktur bei großen Einleitvolumina beträchtlichen Schaden anrichten, da bereits ab einem pH-Wert von unter 5,5 Kleinstlebewesen geschädigt werden \cite{pHAbwasser}.\\
Die Gesamtstickstoffgehalte der Abwässer werden in der Abbildung \ref{nges} dargestellt. Es geht daraus hervor, dass die Probe 1 stark belastet ist. Die Proben 2 und 3 sind im Vergleich zu häuslichen Abwässern sehr gering belastet und erfordern keinen weiteren Handlungsbedarf.\\
Die Abbildung \ref{pges} bildet schließlich die Werte für den Gesamtphosphorgehalt im Vergleich zum häuslichen Abwasser ab. Alle drei untersuchten Proben schneiden in diesem Fall gut, sprich mit einer geringen Belastung, ab.\\

Die Behandlung der untersuchten Abwässer kann durch eine herkömmliche Kläranlage erfolgen. Beim Einleiten in die Kanalisation vermischen sich die Abwässer mit dem Abwasserstrom. Darum könnte man vermuten, dass zum Teil stark belastete Abwässer, wie die Probe 1, verdünnt werden würden und stellen bei Ankunft im Klärwerk keine zusätzliche Herausforderung mehr dar. Der Anteil an der Gesamtabwassermenge würde dabei vernachlässigbar bleiben. 



\chapter{Fehlerbetrachtung}
\label{sec:fehler}



%Praktikumsskript, Modul ………, Versuch …….., Prof. Musterprof. 
%DIN 12345, Jahr der Veröffentlichung 
%Link der Internetseite, Zugriffsdatum 
%Buchtitel, Autor, Verlag, Veröffentlichungsjahr 


%Literaturverzeichnis Bücher
\bibliography{Literatur}
\bibliographystyle{unsrtdin}
\addcontentsline{toc}{chapter}{Literaturverzeichnis}







%Anhang
%\addcontentsline{toc}{chapter}{Anhang}

%\chapter*{Eidesstattliche Erklärung}
\label{erklaerung}
Hiermit versichere ich, die vorliegende Seminararbeit selbstständig und nur unter Verwendung der von mir angegebenen Quellen und Hilfsmittel verfasst zu haben. Sowohl inhaltlich als auch wörtlich entnommene Inhalte wurden als solche kenntlich gemacht. Die Arbeit hat in dieser oder vergleichbarer Form noch keinem anderem Prüfungsgremium vorgelegen. \\
\\[1.5cm]
Datum:	\hrulefill\enspace Unterschrift: \hrulefill
\\[3.5cm]
\addcontentsline{toc}{chapter}{Selbstständigkeitserklärung}

%\chapter{Bausteine}

\section{Beispiel für Tabelle}

%Tabelle START
\vspace*{-2.5mm}
\renewcommand{\arraystretch}{1.2}
\begin{table}[h!]
	\centering
	\caption{Abmessungen der Probekörper vor dem Zugversuch}
	\label{tab:tabelle1}
	%\resizebox{10cm}{!}{
	\begin{tabulary}{\textwidth}{C|CCC}
		\hline
		\textbf{Probe}  &\textbf{Breite [mm]}&\textbf{Dicke [mm]}&\textbf{Anf.-länge[mm]} \\ 
		\hline
		Kupfer (gewalzt) & 12,5 &3,00&50\\
		Kupfer (geglüht) & 12,5&3,00&50\\
		PA6 & 10,0&4,00&50\\
		PP (EPR-30\% Kautschuk) & 9,9 &3,95&50\\
		\hline
	\end{tabulary}
	%}
\end{table}

\FloatBarrier
\vspace*{-2.5mm}
%Tabelle ENDE

\section*{Tabelle mit Itemize}
%Tabelle START
\vspace*{-2.5mm}
\renewcommand{\arraystretch}{1.2}
\begin{table}[h!]
	\centering
	\caption*{Vor- und Nachteile der Geothermie}
	\label{tab:tabelle1}
	\begin{tabulary}{\textwidth}{C|C}
		\hline
		\textbf{Vorteile}  &\textbf{Nachteile} \\ 
		\hline
		&\\
		\begin{minipage}[t]{0.4\textwidth}
			\begin{itemize}
				\item Strom, Wärme und Kälte wird erzeugt
				\item keine saisonalen und tageszeitlichen Schwankungen
				\item 	quasi-regenerativ
				\item 	nachfrage-gerechte Energiebereitstellung
				\item Erzeugungspotenzial sehr hoch 
				\item 	grundsätzlich standortunabhängig
			\end{itemize}
		\end{minipage} & 
		\begin{minipage}[t]{0.4\textwidth}
			\begin{itemize}
				\item hohe Anschaffungskosten
				\item abhängig von geologischen Gegebenheiten
				\item geringer Stromwirkungsgrad (thermodynamisch bedingt)
				\item keine Marktdurchdringung in DE
				\item erfahrene Bauunternehmen notwendig
				\item gute Vorerkundung und Überwachung notwendig
			\end{itemize}
		\end{minipage}\\
	\end{tabulary}
\end{table}
\FloatBarrier
\vspace*{-2.5mm}
%Tabelle ENDE

\newpage

\section*{Beispiel für Skalierbare Tabelle}
%TAbelle Start
\vspace*{-2.5mm}
\renewcommand{\arraystretch}{1.2}
\begin{table}[h!]
	\centering
	\caption*{}
	\resizebox{0.5\textwidth}{!}{
		\begin{tabulary}{\textwidth}{C|C|C|C}
			\textbf{Name} & \textbf{Anwendung}&\textbf{Gleichung}&\textbf{Stoffkonstante} \\ 
			\hline  
			KICK& $x_{80_\omega}>\SI{50}{\milli\meter}$ &$e_{KICK}=c_K*log(\frac{x_{80_\omega}}{x_{80_\alpha}})$&$c_K=1,15*\frac{c_B}{\sqrt{0,05\si{\meter}}} \left[\si{\raiseto{2}\meter\per\raiseto{2}\second}\right]$\\
			BOND&$\SI{50}{\micro\meter}<x_{80_\omega}<\SI{50}{\milli\meter}$&$e_{BOND}=c_B*\left(\frac{1}{\sqrt{x_{80_\omega}}}-\frac{1}{\sqrt{x_{80_\alpha}}}\right)$& $c_B$: tabelliert $\left[\si{\raiseto{2,5}\meter\per\raiseto{2}\second}\right]$\\ 
			RITTER& $x_{80_\omega}>\SI{50}{\micro\meter}$&$e_{RITT}=c_R*\left(\frac{1}{x_{80_\omega}}-\frac{1}{x_{80_\alpha}}\right)$&$c_R= 0,5*c_B*\sqrt{\SI{5e-5}{\meter}}$ \\  
	\end{tabulary}}
\end{table}
\FloatBarrier
%Ende TAbelle

\section{Befehlszeilen in Text einfügen}
Befehlzeilen aus der "Latex-Sprache"\ lassen sich nicht ohne weiteres im Text darstellen. Das System erkennt diese als solche und gibt warnhinweise aus. In der Regel werden die Befehle dann auch falsch dargestellt. Umgehen lässt sich diese Problematik mit der verbatim-Umgebung. In ihren Grenzen werden Eingaben 1:1 dargestellt wie eingegeben. Die Funktionsweise soll nachfolgend durch ein Beispiel verdeutlicht werden.
\begin{verbatim*}

\begin{verbatim}
\usepackage{Beispieltrolllllollllolll} 
\end{verbatim}

\end{verbatim*}

Sollte besonderer Wert auf die Kenntlichmachung der Leerzeichen gelegt werden kann auch mit \texttt{verbatim*} gearbeitet werden.

\section{Seitenübergreifende, lange Tabellen}

Tabellen welche Messwerte in einem solchen Umfang enthalten, dass sie nicht auf einer einzelnen A4-Seite Platz finden, können mit dem Paket \begin{verbatim}
\usepackage{longtable} 
\end{verbatim}
in ein Dokument eingepflegt werden wie folgendes Beispiel belegt.:
\begin{longtable}[c]{lllll}
	\caption{Dehnungstabelle}\\
	\label{alles}
	$Zeit [HH:MM:SS]$ & $\Delta l_{PE} [mm]$ & $\varepsilon_{PE}$ & $\Delta l_{Pb} [mm]$ & $\varepsilon_{Pb}$ \\
	\hline
	\endfirsthead
	%\caption{}\\
	$Zeit [HH:MM:SS]$ & $\Delta l_{PE} [mm]$ & $\varepsilon_{PE}$ & $\Delta l_{Pb} [mm]$ & $\varepsilon_{Pb}$ \\ 
	\hline
	\endhead
	\multicolumn{5}{r}{Fortsetzung auf n{\"a}chster Seite}\\
	\endfoot
	\hline
	\multicolumn{5}{r}{} \\
	\endlastfoot
	% Ab hier kommt der Inhalt der Tabelle
	00:00:00 & 3,00 & 0,11 & 3,30 & 0,12\\
	00:00:10 & 3,80 & 0,14 & 3,42 & 0,13\\
	00:00:20 & 4,15 & 0,16 & 3,53 & 0,13\\
	00:00:30 & 4,43 & 0,17 & 3,68 & 0,14\\
	00:00:40 & 4,61 & 0,17 & 3,71 & 0,14\\
	00:00:50 & 4,80 & 0,18 & 3,81 & 0,14\\
	00:01:00 & 4,95 & 0,19 & 4,12 & 0,15\\
	00:20:00 & 8,03 & 0,30 &  & \\
\end{longtable} 

\section{Diagramme}
Diagramme lassen sich unter anderem mit dem Paket 
\begin{verbatim}
\usepackage{pgfplots} 
\end{verbatim}
implementieren.
Um das Dokument einigermaßen klein zu halten empfehle ich folgendes weiters Paket zu nutzen.
\begin{verbatim}
\usepackage{csvsimple} 
\end{verbatim}
Es erleichter das importieren von Datensätzen aus .csv Dateien. Die .csv Datei muss folgenden Anforderungen genügen:- Dezimaltrenner Punkt
\begin{itemize}
	\item Dezimaltrenner Punkt
	\item Jede Zeile umgebrochen
	\item Koordinaten durch komma getrennt
	\item Spalten beschriften z.B x,y oder a,b
\end{itemize}

Sollte man eine Tabellenkalkulationsdatei in eine csv Datei konvertiert haben, hat dieses meist die Falsche Form. Durch die Funktion suchen und ersetzen (bsp emacs) können aber sehr schnell die nötigen Korrekturen erfolgen. 

\begin{figure}[h]
	\begin{center}
		\begin{tikzpicture}
		\begin{axis}[
		width=12cm,
		height=6cm,
		xlabel=Zeit in Sekunden,
		ylabel=Dehnung]
		
		\begin{scope}[brown]
		\draw[brown] ({axis cs:10,0}|-{rel axis cs:0,1}) -- ({axis cs:10,0}|-{rel axis cs:0,0});
		\draw[brown] ({axis cs:120,0}|-{rel axis cs:0,1}) -- ({axis cs:120,0}|-{rel axis cs:0,0});
		\end{scope} 
		
		\addplot table [x=a, y=b, col sep=comma] {data/KriechkurvePb.csv};
		\end{axis}
		\end{tikzpicture}
		\caption{Kriechkurve Blei}
		\label{kkb}
	\end{center}
\end{figure} 
\FloatBarrier                     

\section*{Beispiel für Berechnungen}
Die Berechnung der wahren Spannung $\sigma_{W}$ bei Höchstkraft erfolgt unter der Annahme, dass der Prüfkörperquerschnitt noch 60\% des Ausgangsquerschnitts beträgt. Die Berechnung erfolgt ab Gleichung \ref{ber1}.\\ 

%Berechnung der Fläche A
\textbf{Fläche} $\boldsymbol{A}$ \textbf{:}
%Start
\begin{flalign}
A 	&= \frac{\pi}{4}*d^2\\
&=\frac{\pi}{4}*(\SI{80}{\milli \meter})^2
\end{flalign}
%Ende

\section{Beispiel für ein Bild}

%Start
\begin{figure}[h!]
	\centering
	\includegraphics[width=0.60\textwidth]{img/skizzepruef3}
	\caption{Skizze Prüfkörperbemaßung}
	\label{skizzepruef}
\end{figure}
\FloatBarrier
%Ende

\newpage

\section{Beispiel für zwei Bilder}
\label{sec:versuchsaufbau}

%Start
\begin{figure}[h!]
	\centering
	\begin{subfigure}{.5\textwidth}
		\centering
		\includegraphics[width=0.75\textwidth]{Aufbau2}
		\caption{Skizze zum Versuchsaufbau}
		\label{fig:sub1}
	\end{subfigure}%
	\begin{subfigure}{.5\textwidth}
		\centering
		\includegraphics[width=0.6\textwidth]{img/Aufbau1}
		\caption{realer Versuchsaufbau}
		\label{fig:sub2}
	\end{subfigure}
	\caption{Versuchsaufbau als Skizze und in Realität}
	\label{fig:aufbau} 
\end{figure}
\FloatBarrier
%Ende

\newpage


\section{Beispiel für vier Bilder}
%Start
\begin{figure}[h!]
	\centering
	\begin{subfigure}{.5\textwidth}
		\centering
		\includegraphics[width=0.75\textwidth]{img/Kupfer_gewalzt}
		\caption{Kupfer (gewalzt)}
		\label{fig:sub3}
	\end{subfigure}%
	\begin{subfigure}{.5\textwidth}
		\centering
		\includegraphics[width=0.75\textwidth]{img/Kupfer_weich}
		\caption{Kupfer (geglüht)}
		\label{fig:sub4}
	\end{subfigure}
	\begin{subfigure}{.5\textwidth}
		\centering
		\includegraphics[width=0.75\textwidth]{img/PA6}
		\caption{PA6}
		\label{fig:sub5}
	\end{subfigure}%
	\begin{subfigure}{.5\textwidth}
		\centering
		\includegraphics[width=0.75\textwidth]{img/PP}
		\caption{PP (ERP-30\% Kautschuk)}
		\label{fig:sub6}
	\end{subfigure}
	
	\caption{Bruchstellennahaufnahmen der Probekörper}
	\label{fig:bruchstellen} 
\end{figure}
\FloatBarrier
%Ende

\section{Beispiel Einheiten}

\begin{align*}
\textbf{$\SI{12,0/12}{\kg\meter\per\second \raiseto{5} \per \xyz}*\SI{13}{\per\raiseto{-2}\meter}=\SI{256}{}$}
\end{align*}

\begin{align}
\SI{12,0/12}{\meter\per\joule}*\SI{13}{\gram}=\SI{256}{\coulomb}
\end{align}

\newpage

\section{Beispiel für Mini-Formelsammlung}
\begin{flalign}
\label{gl1}
\text{\textbf{Dehnung (Def.)} } \boldsymbol{\varepsilon} \text{ \textbf{:}} && \hspace*{-1em}  \varepsilon=\frac{\Delta l}{l_0} &&
\end{flalign}

\begin{flalign}
\label{gl2}
\text{\textbf{norminelle Spannung} } \boldsymbol{\sigma} \text{\textbf{:}} && \hspace*{-3em} \sigma=\frac{F}{A_0} &&
\end{flalign}

\begin{flalign}
\label{gl3}
\text{\textbf{Sekantenmodul (Kunststoffe) }} \boldsymbol{E_S} \text{ \textbf{:}} && E_S=\frac{\sigma_2-\sigma_1}{\varepsilon_2-\varepsilon_1}=\frac{F_2-F_1}{0,002*A_0} &&
\end{flalign}

\begin{flalign}
\label{gl4}
\text{\textbf{E-Modul (Metalle) }} \boldsymbol{E_M} \text{ \textbf{:}} && \hspace*{5em} E_M=\frac{\sigma}{\varepsilon}=\frac{\sigma_2-\sigma_1}{\varepsilon_2-\varepsilon_1}=\frac{R_{p_{0.2\%}}}{0,2\%} &&
\end{flalign}

\begin{flalign}
\label{gl5}
\text{\textbf{Bruchdehnung} } \boldsymbol{A} \text{\textbf{:}} && \hspace*{6em} A= \frac{l_{u}-l_{0}}{l_{0}}*100\% &&
\end{flalign}

\begin{flalign}
\label{gl6}
\text{\textbf{Ausgangsquerschnitt } } \boldsymbol{S_0} \text{\textbf{:}} && \hspace*{3em} S_{0}= Breite*Dicke &&
\end{flalign}
\begin{flalign}
\label{gl7}
\text{\textbf{wahre Spannung} } \boldsymbol{\sigma_{W}} \text{\textbf{:}} &&\hspace*{1em} \sigma_{W}=\frac{F_{max}}{S_{End}}&&
\end{flalign}
\begin{flalign}
\label{gl8}
\text{\textbf{Brucheinschnürung } } \boldsymbol{Z} \text{\textbf{:}} && \hspace*{5em} Z=\frac{S_0-S_u}{S_{o}}*100\% &&
\end{flalign}

\newpage

\section*{Fußnoten}
%Start
\begin{figure}[h!]
	\centering
	\includegraphics[width=0.85\textwidth]{tabdia/kfwerte}
	\caption*{$\text{k}_\text{f}$-Werte der Proben 1 bis 11 \protect\footnotemark[1]}
	\label{}
\end{figure}
\FloatBarrier
%Ende

\footnotetext[1]{bezogen auf $V=\SI{50}{\milli \liter}$ und $h=\SI{10}{\milli \meter}$}



\end{document}
